% ---------------------------------------------------------------------------
% riccati-cert-simple/main.tex -- the SIMPLIFIED paper (the publication target).
% Stage 4 skeleton + prose. Seat 1 wrote: preamble, Abstract, \S1, \S2 (incl.
% \S2.1). Sections 3--9 are placeholders for the later seats.
%
% LABEL SCHEME (binding for all seats). Section labels REUSE the full paper's
% names, because the Stage-3 plan-map's claim texts embed them (\S\ref{sec:rpod}
% and so on); the simplified-only additions follow the same prefixes:
%   sections:    sec:intro  sec:setup  sec:sizes (=\S2.1)  sec:riccati
%                sec:theorem  sec:regime  sec:rpod  sec:related  sec:limits
%                sec:conclusion (new; the full paper's sec:limits covers both)
%   equations:   eq:qcqp  eq:keepout  eq:riccati
%   theorem envs: def:height  def:protocol  def:constructed  thm:main
%                lem:band  lem:dd  ex:gap  rem:value  rem:effective
%                rem:curvature (new)  rem:bordered (new)  rem:contact (new)
%   tables:      tab:heights  tab:results
% New numbered objects introduced by the restructure take the same prefixes
% (def:, lem:, rem:, ex:, tab:, eq:).
% ---------------------------------------------------------------------------
\documentclass[11pt]{article}
\usepackage[margin=1in]{geometry}
\usepackage{amsmath,amssymb,amsthm}
\usepackage{booktabs}
\usepackage[hidelinks]{hyperref}

\newtheorem{theorem}{Theorem}
\newtheorem{lemma}{Lemma}
\theoremstyle{definition}
\newtheorem{assumption}{Assumption}
\newtheorem{definition}{Definition}
\newtheorem{example}{Example}
\theoremstyle{remark}
\newtheorem{remark}{Remark}
\newcommand{\Q}{\mathbb{Q}}
\newcommand{\R}{\mathbb{R}}
\DeclareMathOperator{\lcm}{lcm}
\DeclareMathOperator{\den}{den}

% Typography: long typewriter identifiers carry no break points, so a line
% ending in one overflows the measure. Allow a break after each underscore,
% and give TeX a third-pass stretch budget so an unbreakable token moves to
% the next line instead of running past it. check_refs fails the build on any
% overfull box wider than 2pt, so this must hold.
% OT1's \_ is a drawn RULE, not a character, so identifiers extracted from the
% PDF lost their underscores -- copy-paste and any text indexer saw the wrong
% name. Emit the typewriter font's real underscore glyph and let pdftex map it
% to U+005F. The glyph exists only in the typewriter font -- char 95 of cmr is
% not an underscore -- so the roman case (QSopt\_ex) keeps the rule. The
% fontdimen3 test is the standard "is this monospace" check: tt fonts have no
% interword stretch.
\input{glyphtounicode}
\pdfgentounicode=1
\renewcommand{\_}{\ifdim\fontdimen3\font=0pt\relax\char`\_\else\textunderscore\fi\hspace{0pt}}
\emergencystretch=3em

% Draft version. SINGLE SOURCE OF TRUTH for the version number: papers/build.sh
% names the output <paper-dir>-v<paperversion>.pdf from this macro, and each
% paper's verify/_common.py parses it to find the build products. Bump it here
% and nowhere else; the title page carries it so a loose PDF names its version.
\newcommand{\paperversion}{1}

\title{Certifying Trajectory Optimality in Exact Arithmetic:\\
An Exact-Rational Riccati Certificate for Penalty-Coupled Trajectory QCQPs}
\author{Adi Oltean\thanks{Starcloud, Inc. \texttt{adi@starcloud.com}. Working draft.
AI coding assistants were used for drafting, for the reference implementation and
verification scripts, and for adversarial review; all results were checked by the author.}}
\date{Version \paperversion{} --- August 2026}

\begin{document}
\maketitle

\begin{abstract}
Trajectory optimization for spacecraft rendezvous under nonconvex keep-out
(\emph{state}) constraints yields nonconvex quadratically constrained quadratic
programs (QCQPs) whose global optimality is, in practice, certified only
numerically, up to a solver tolerance; we treat the \emph{penalty-coupled}
(soft-dynamics) subclass, in which the stage-to-stage dynamics enter the cost
rather than as equality constraints and the decision variables are the stage
states alone, with no control variable. For this class, the S-procedure
positive-semidefiniteness (PSD) test restricted to the problem's band \emph{is}
the discrete LQ Riccati (block-Schur) recursion, so running it in exact rational
arithmetic gives a tolerance-free certified lower bound in $O(Nd^3)$ exact
rational operations; the band identity is proved, while the border-carrying
extension that delivers the scalar corner has its soundness asserted and
test-pinned, its count
read off the sweep's loop bounds. Under an aggregate-PSD condition at a
constructed rational optimum the bracket closes exactly, upgrading the bound to
a global-optimality certificate. The construction places the emitted certificate
$(\lambda^\star,x^\star)$ at per-stage design height $L$ before any datum is
derived from it, so the emitted pair has per-entry bit-size $O(L)$ independent
of the horizon; the theorem's content is that this designed pair is the
certified global optimum and that the verifier recovers it at that size. The
certified value the verifier recomputes is not bounded by $L$, though: it is
governed instead by the instance's assembled height together with the height of
the placed coordinates and a $\log N$ term, with neither ordering between the
assembled and design heights forced. No unconditional bound exists for quadratic
optimization at large, which admits $2^n$-bit exact optima (whether such
instances live in our class is open), so the guarantee is
stated conditionally. Aggregate-PSD is implied by a solver-free
weak-diagonal-dominance rule, and a counterexample certified in exact
arithmetic shows it is a real assumption. The trusted core is machine-checked in
Lean~\cite{leannote}, and the pipeline is demonstrated end-to-end on an
inverse-designed Clohessy--Wiltshire-coupled keep-out stress instance, checked
throughout by an exact implementation that refuses floats. The protocol is
sound but not complete: acceptance is a theorem about the deployed numbers,
and refusal asserts nothing.
\end{abstract}

\section{Introduction}\label{sec:intro}
Guidance for spacecraft rendezvous, proximity operations, and docking (RPOD)
plans a trajectory by solving optimal-control problems that, after
discretization, are QCQPs: a quadratic cost (tracking, path smoothness) subject
to quadratic keep-out constraints such as obstacles and abort-safety balls.
Corridors and other convex tubes carry the opposite sign and are not instances
of the class treated here. We treat throughout the \emph{penalty-coupled}
(soft-dynamics) subclass, in which the stage-to-stage dynamics enter the cost
rather than as equality constraints and the decision variables are the stage
states alone, with no control variable; hard equality dynamics are outside the
class treated here (\S\ref{sec:limits}).

These QCQPs are nonconvex, and standard practice --- sequential convex
programming, or numerical moment--SOS
relaxations~\cite{lasserre2001,strom2024,yang2023perception} --- returns a
solution whose global optimality is either not certified at all or certified in
floating point, up to a tolerance. Lossless
convexification~\cite{acikmese2011} does certify exactly and analytically, via
the maximum principle, but for nonconvex \emph{control}-set constraints --- not
the state constraints at issue here.

We ask a different question: can the \emph{certificate} of global optimality of
a trajectory QCQP be made \emph{exact}, checked in rational arithmetic with no
tolerance, so that acceptance is a theorem about the deployed numbers rather
than a numerical near-miss? For a single keep-out the answer is known --- the
S-lemma is exact and its dual is one PSD
test~\cite{polik2007,yakubovich1971}, with rational
recovery~\cite{peyrl2008,kaltofen2012}; the open case is the \emph{trajectory},
a chain of coupled primitives, where exact certification might fail (a duality
gap) or explode (exponential-bit rationals~\cite{bienstock2020}).
The question is about checking, not searching. A proposer supplies the
candidate --- a numerical solver in general, the inverse design of
\S\ref{sec:rpod} in our demonstrations --- and nothing in the certificate
route replaces the search that produced it.

\paragraph{The structural observation.} The S-procedure certificate for the
chain is one bordered PSD condition $M(\lambda,t)\succeq0$ (\S\ref{sec:setup}).
Its exact test is a rational $LDL^\top$ that eliminates within the problem's
band and carries the affine column alongside; that test is precisely the
discrete Riccati / block-Schur recursion of LQ optimal-control
numerics~\cite{wright1993,rao1998}, and the recursion's terminal arrival cost
delivers the single scalar corner condition. Run in exact $\Q$, acceptance
becomes a tolerance-free theorem at the $O(Nd^3)$ operation count of one LQ
Riccati sweep, modulo the border-carrying extension of \S\ref{sec:riccati}.

\paragraph{Delta over the published technical note.} The published Podium
technical note~\cite{oltean2026certnote} establishes the exact S-procedure
bracket and the rational $LDL^\top$ PSD test for \emph{single} primitives; the
mathematics there is classical and presented as such. This paper's delta is
the \emph{trajectory} case:
\begin{itemize}
\item the identity between the band-restricted test and the Riccati recursion,
carried out in $\Q$;
\item conditional exact closure with a horizon-independent emitted certificate
(Theorem~\ref{thm:main});
\item the solver-free diagonal-dominance criterion, with a gap counterexample
certified in exact arithmetic (\S\ref{sec:regime});
\item band checkers shipped in the Podium library; and
\item a machine-checked Lean core~\cite{leannote}.
\end{itemize}

\paragraph{Contributions.}
\begin{enumerate}
\item \textbf{Exact-$\Q$ soundness (\S\ref{sec:riccati}).} The band-restricted
rational $LDL^\top$ equals the dense PSD test, with exact zero-pivot handling,
in $O(Nd^3)$ rational operations (Lemma~\ref{lem:band}, \emph{proved}). The
\emph{acceptance predicate} --- a nonnegative-pivot factorization exists iff
the matrix is PSD --- is \emph{machine-checked} in Lean in both directions;
that the shipped sweep computes such a factorization is not machine-checked
(\S\ref{sec:riccati}, \cite{leannote}). The bordered sweep ships as a trusted
checker in the Podium library.
\item \textbf{Conditional closure (\S\ref{sec:theorem}).} Under aggregate-PSD
(A1), a rational recoverable optimum (A2), and linear independence of the
active constraint gradients (LICQ), the bracket closes exactly, with any
number of keep-outs per stage; LICQ leaves at most $d$ active on each stage.
For a \emph{penalty-coupled} constructed rational-sphere family
(Definition~\ref{def:constructed}) the emitted certificate
$(\lambda^\star,x^\star)$ has \emph{per-entry} bit-size $O(L)$ independent of
$N$, by the design order of Definition~\ref{def:constructed} rather than by a
measurement. The theorem's content is that the pair so placed is the certified
global optimum and is recovered by the verifier at that size; both legs are
\emph{proved}. The guarantee is conditional, and \S\ref{sec:theorem} says why
no unconditional one is available.
\item \textbf{A solver-free benign criterion (\S\ref{sec:regime}).} Weak
diagonal dominance implies aggregate-PSD and is preserved along the sweep
(Lemma~\ref{lem:dd}, \emph{proved}; core
\emph{machine-checked}~\cite{leannote}). A duality-gap counterexample
certified in exact arithmetic (Example~\ref{ex:gap}) shows aggregate-PSD is a
real assumption, not automatic. The certification pipeline is \emph{demonstrated}
end-to-end on an inverse-designed CW-coupled stress instance
(\S\ref{sec:rpod}), where aggregate-PSD is verified by the general sweep
rather than by this solver-free shortcut; every reported number is reproduced
by an executable exact-arithmetic receipt (\S\ref{sec:rpod},
\cite{leannote}).
\end{enumerate}

\section{Problem and certificate protocol}\label{sec:setup}
Let $x=(x_1,\dots,x_N)$ with each $x_k\in\R^d$ and total dimension $n=Nd$; the
optimization is over $\R$, while the problem \emph{data} and the
\emph{certificate} are rational. Consider
\begin{equation}\label{eq:qcqp}
J^\star=\min_x\; f_0(x)\quad\text{s.t.}\quad f_j(x)\ge 0,\ j=1,\dots,m,
\end{equation}
with a block-tridiagonal objective coupling only neighbours,
\[
f_0(x)=\sum_{k=1}^{N-1} \lVert x_{k+1}-A_k x_k\rVert_{R_k}^2
       +\sum_{k=1}^{N} (x_k-g_k)^\top W_k (x_k-g_k),
\]
with $R_k\succ0$ the stage-cost weight and $W_k$ the goal weight, and per-stage
keep-outs
\begin{equation}\label{eq:keepout}
f_j(x)=\lVert \Pi_j x_{\sigma(j)}-c_j\rVert^2-\rho_j^2,\qquad \rho_j>0,
\end{equation}
each involving only the single stage $\sigma(j)\in\{1,\dots,N\}$; several
keep-outs may guard the same stage. Here $\Pi_j$ is a fixed coordinate
projection selecting the \emph{guarded} sub-block of the stage state:
$\Pi_j=I$ is the full-stage keep-out ball, while a guidance keep-out typically
guards the position sub-block of a position--velocity state (\S\ref{sec:rpod}).
The neighbour-only coupling confines every nonzero entry to a scalar band of
half-width $b=2d-1$, and that is the band the sweep of \S\ref{sec:riccati}
uses.

The certificate's shape can be read off the one-dimensional picture first.
At the minimum of a parabola the tangent is horizontal: the derivative
vanishes at the minimizer, and that first-order identity is an algebraic
fact about the minimizer's own coordinates --- a claimed minimizer is
checked by evaluating a derivative, not by rerunning any solver.
When the minimum sits on the boundary of the feasible set ---
here, on an active keep-out sphere --- the gradient need not vanish.
Under a qualification on the active gradients (the theorem's LICQ), it is
balanced by the active constraints' gradients with nonnegative weights, and
those weights are the multipliers $\lambda$ of the KKT stationarity
condition. The balance alone does not certify. With nonconvex keep-outs
many points satisfy it: local minima, saddle points, contacts on the wrong
branch of a sphere. And which keep-outs are active at the optimum is a
combinatorial unknown. The S-procedure certificate $(\lambda,t)$ defined
next is the object that closes the gap: acceptance upgrades ``these
coordinates balance'' to the global statement $t\le J^\star$, and closure
of the bracket (Definition~\ref{def:protocol}) upgrades it again, to global
optimality of the claimed point --- checked in exact arithmetic throughout.

Write $f_0(x)=x^\top P_0 x+q_0^\top x+r_0$ with $P_0$ block-tridiagonal, and
$f_j(x)=x^\top P_j x+q_j^\top x+r_j$. Then $P_j=\Pi_j^\top\Pi_j\succeq0$ is
supported on the single block $\sigma(j)$, and since $\Pi_j$ is a coordinate
projection, $\lVert\nabla f_j\rVert=2\rho_j$ on the contact sphere. These are
the structural properties of the keep-outs that the results below use;
Lemma~\ref{lem:dd}'s sharp form additionally takes $A_k=I$, $R_k=I$, $W_k=w_kI$
and $\Pi_j=I$. A multiplier vector $\lambda\in\Q_{\ge0}^m$ and a value $t\in\Q$
certify the lower bound $t\le J^\star$ when
\[
M(\lambda,t)=\begin{bmatrix} H(\lambda) & \tfrac12(q_0-\textstyle\sum_j\lambda_j q_j)\\[2pt]
\tfrac12(q_0-\sum_j\lambda_j q_j)^\top & r_0-\sum_j\lambda_j r_j-t\end{bmatrix}\succeq 0,
\qquad H(\lambda):=P_0-\textstyle\sum_j\lambda_j P_j,
\]
by weak duality: $M\succeq0$ makes $f_0-\sum_j\lambda_j f_j-t$ globally
nonnegative, so $t+\sum_j\lambda_j f_j$ is a global minorant of $f_0$. The
converse fails: Example~\ref{ex:gap} exhibits true lower bounds admitting no
level-one certificate. Acceptance of a certificate
(Definition~\ref{def:protocol}) is therefore the sufficient direction.

By a generalized Schur complement, $M(\lambda,t)\succeq0$ holds iff three
conditions do:
\begin{enumerate}
\item[(i)] $H(\lambda)\succeq0$;
\item[(ii)] the linear part $\tfrac12(q_0-\sum_j\lambda_j q_j)$ lies in
$\mathrm{range}(H(\lambda))$ --- automatic when $H$ is nonsingular, and when
$H$ is singular the sweep detects the condition as a consistent zero-pivot
row;
\item[(iii)] the scalar corner condition $t\le g(\lambda)$, where $g(\lambda)$
is the Lagrangian dual function value at $\lambda$.
\end{enumerate}

\begin{remark}[Curvature convention]\label{rem:curvature}
Block tridiagonality of $H(\lambda)$ is forced by the class rather than
assumed: every summand of $f_0$ touches only $\{x_k,x_{k+1}\}$ or $x_k$ alone,
each $P_j$ sits in the single block $\sigma(j)$, and block-tridiagonal
matrices form a linear subspace. The structure is load-bearing: the band sweep
cannot see out-of-band entries, and it returns a wrong verdict on a symmetric
matrix that lacks the structure. Under the convention
$f_0=x^\top P_0x+q_0^\top x+r_0$ used here, with no $\tfrac12$ factor,
$H(\lambda)$ is half the Lagrangian Hessian:
$\nabla^2_{xx}(f_0-\sum_j\lambda_jf_j)=2H(\lambda)$. We therefore call
$H(\lambda)$ the \emph{Lagrangian curvature matrix}.
The factor is fixed by that convention
and is load-bearing wherever $H$ enters quantitatively: $H$, and not $2H$, is
the curvature block of $M(\lambda,t)$, and the pivot bit-sizes
\S\ref{sec:rpod} measures are those of $H$. What is invariant under a positive
rescaling are the three verdict-level properties invoked below ---
semidefiniteness, weak diagonal dominance, and pivot signs --- since the band
pivots are homogeneous of degree one in $H$.
\end{remark}

\begin{definition}[Certificate protocol]\label{def:protocol}
The \emph{S-procedure certificate} is the pair $(\lambda,t)$; its acceptance
--- $\lambda\ge0$ and $M(\lambda,t)\succeq0$, checked in exact arithmetic ---
proves $t\le J^\star$. The \emph{emitted certificate} is the object a prover
transmits: the pair $(\lambda,\bar x)$ of multipliers and a witness. The
verifier checks $\lambda\ge0$ and $f_j(\bar x)\ge0$ for every $j$.
It then recomputes both bounds: a candidate $t$ as the Lagrangian
value at $\bar x$, which is the route the shipped drivers take, with
acceptance re-verifying $M(\lambda,t)\succeq0$ however $t$ was formed; and the
upper bound as $f_0(\bar x)$. The certified value is therefore never
transmitted or trusted.
The two outcomes are asymmetric, in the standard sense that the protocol is
\emph{sound} but not complete. Acceptance carries the bound as a theorem: no
tolerance exists through which a false certificate could pass. A candidate
failing any check is \emph{refused}, and a refusal asserts nothing --- in
particular, it does not prove the candidate suboptimal (\S\ref{sec:regime}
exhibits an instance whose true optimum admits no level-one certificate at
all).
Throughout, the verifier works in exact rational arithmetic, evaluating $f_0$
and the $f_j$ and assembling $M(\lambda,t)$ from its own copy of the problem
data: a certificate is meaningful only relative to agreed data. The interval $[t,f_0(\bar x)]$
\emph{brackets} $J^\star$, and the bracket \emph{closes} when $t=f_0(\bar x)$.
Every size or cost claim below names which object it concerns; for claims
about a certificate, that is the S-procedure certificate or the emitted
certificate.
\end{definition}

\subsection{Size measures}\label{sec:sizes}
All data are rational, and every bit bound in this paper is stated against a
height --- in the sense of Definition~\ref{def:height} --- of a tuple that the
bound names. Two statements stand outside the scheme: the per-entry bit-size of
the rationalized Clohessy--Wiltshire (CW) state-transition matrix in
\S\ref{sec:rpod}, stated against the denominator cap that produces it, and the
imported $2^n$-bit lower bound of~\cite{bienstock2020}, stated against the
problem dimension.

\begin{definition}[Height and bit-size]\label{def:height}\leavevmode
\begin{enumerate}
\item[(a)] For a finite tuple $T$ of rationals, its \emph{height}
$\mathrm{ht}(T)$ is the least $h$ such that $T$ can be written over one common
denominator, $T=(a_1/D,\dots,a_r/D)$ with $D$ a positive integer and the $a_i$
integers, subject to $\lceil\log_2(D{+}1)\rceil\le h$ and
$\lceil\log_2(|a_i|{+}1)\rceil\le h$. That least $h$ is attained at the least
common multiple of the entries' reduced denominators, which is the $D$ the
receipts compute. Clearing that denominator turns the tuple into integers of
at most $h$ bits --- which is what Hadamard's and Edmonds' bounds require.
\item[(b)] A single rational $p/q$ in lowest terms has bit-size
$\mathrm{bit}(p/q)=\lceil\log_2(|p|{+}1)\rceil+\lceil\log_2(q{+}1)\rceil\le
2\,\mathrm{ht}(p/q)$; that per-entry measure is what Table~\ref{tab:results}
reports for the emitted certificate.
\item[(c)] For a positive \emph{integer} $D$ we write
$\mathrm{bitlen}(D)=\lceil\log_2(D{+}1)\rceil$ for its bit-length, the measure
the receipts apply to denominators; read as a rational,
$\mathrm{bit}(D)=\mathrm{bitlen}(D)+1$.
\end{enumerate}
\end{definition}

\noindent Table~\ref{tab:heights} collects the named heights the bounds below
are stated against; each is the height, in the sense of
Definition~\ref{def:height}, of a different tuple, and \S\ref{sec:rpod}
measures all four.

\begin{table}[ht]\centering
\caption{The named heights.}
\label{tab:heights}
\begin{tabular*}{6.1in}{@{}p{1.35in}p{2.35in}p{2.0in}@{}}
\toprule
Height & tuple measured & where used\\
\midrule
$L$ (\emph{per-stage design height}) & the design tuple
Definition~\ref{def:constructed} places on a stage: that stage's placed
coordinates, its active keep-out centres $c_j$, and its active multipliers
$\lambda^\star_j$ & Theorem~\ref{thm:main}'s bit-size leg, stated against this
tuple and no larger one\\
\addlinespace
$L_{\mathrm{asm}}$ (assembled height) & the assembled coefficients
$P_0,q_0,r_0$, and the goals $g_k$ they are recovered from &
Remark~\ref{rem:value}'s bounds on the certified value\\
\addlinespace
$L_M$ (bordered height) & the bordered matrix $M(\lambda,t)$ that the sweep
eliminates: $L_M=\mathrm{ht}(M(\lambda,t))$ & Remark~\ref{rem:effective}'s
pivot bit-size bound\\
\addlinespace
$\mathrm{ht}(x^\star)$ & the placed coordinates taken over \emph{all} stages
at once & Remark~\ref{rem:value}'s sharper first inequality\\
\bottomrule
\end{tabular*}
\end{table}

\noindent These are different tuples and in general different numbers, and no
ordering among $L$, $L_{\mathrm{asm}}$, $L_M$ and $\mathrm{ht}(x^\star)$ is
forced, so each bound below names which height it uses. The corner
$r_0-\sum_j\lambda_jr_j-t$ of the bordered matrix can cancel $r_0$'s
denominator against $t$'s, and on the demonstration family of \S\ref{sec:rpod}
it does. Nor is $\mathrm{ht}(x^\star)$ a renaming of $L$: $L$ bounds a stage's
design tuple one stage at a time, while $\mathrm{ht}(x^\star)$ is taken over
the placed coordinates of all stages at once.

\section{The certificate is a Riccati sweep}\label{sec:riccati}
$H=H(\lambda)$ is symmetric block-tridiagonal, with $d\times d$ blocks
$H_{kk}$ and couplings $H_{k,k+1}=-A_k^\top R_k$. Its scalar bandwidth is at
most $b=2d-1$. Equality holds whenever some coupling block has a nonzero
$(1,d)$ entry, as it does on the family of \S\ref{sec:rpod}; structured
couplings such as $A_k=R_k=I$ instead sit at half-width exactly $d$, strictly
inside the band whenever $d\ge2$. The sweep uses the conservative $b$ in
either case. Eliminating blocks in increasing order
produces the pivots
\begin{equation}\label{eq:riccati}
S_1=H_{11},\qquad S_k=H_{kk}-H_{k,k-1}\,S_{k-1}^{-1}\,H_{k-1,k}
\quad(\text{while } S_{k-1}\succ0),
\end{equation}
and, in the positive-definite interior, $H\succ0$ iff every $S_k\succ0$; the
general semidefinite verdict, including singular boundary pivots, is delivered
by the scalar band-$LDL^\top$ of Lemma~\ref{lem:band} with its zero-pivot
handling.

Equation~\eqref{eq:riccati} is the discrete \emph{Riccati / block-Schur}
recursion of an LQ problem: $S_k$ is the accumulated (arrival-cost /
information-form) curvature matrix after eliminating $x_1,\dots,x_{k-1}$ --- a
Schur complement in $H$, so it inherits the same no-$\tfrac12$ convention as
$H(\lambda)$ (Remark~\ref{rem:curvature}). That the block-banded KKT
factorization equals the Riccati recursion is
classical~\cite{wright1993,rao1998}; here it is run in exact $\Q$.

The full certificate is $M(\lambda,t)\succeq0$, that is, the three conditions
of \S\ref{sec:setup}. Carrying the affine term
$\tfrac12(q_0-\sum_j\lambda_j q_j)$ through the same sweep makes the terminal
corner pivot $g(\lambda)-t$; equivalently, run the sweep at $t=0$ and read off
the arrival cost $g(\lambda)$. All three conditions are thus certified in one
tolerance-free rational sweep of $O(Nd^3)$ operations. That sweep is an
extension of Lemma~\ref{lem:band}; its \emph{soundness} is asserted and
test-pinned rather than proved, and its operation count is read off the
sweep's loop bounds (Remark~\ref{rem:bordered}).

\begin{lemma}[Band-$LDL^\top$ soundness]\label{lem:band}
For symmetric $H$ with bandwidth $b$, the $LDL^\top$ that eliminates only
within the band yields the same pivots --- hence the same PSD verdict --- as
the dense test, in $O(nb^2)=O(Nd^3)$ exact rational operations, with exact
zero-pivot handling.
\end{lemma}
\noindent Band preservation is proved below; the content of the lemma is the
exact-rational, zero-pivot-handling PSD \emph{verdict}.
\begin{proof}
Eliminating pivot $k$ updates $a_{ij}\mathrel{-}=(a_{ik}/a_{kk})a_{kj}$ for
$i,j>k$. With bandwidth $b$, $a_{ik}\ne0\Rightarrow i\le k+b$ and
$a_{kj}\ne0\Rightarrow j\le k+b$, so only entries with $i,j\in(k,k+b]$ change
and the bandwidth is preserved; by induction all nonzero updates lie in the
band, so eliminating only there is exact. At a zero pivot, positive
semidefiniteness requires $a_{kj}=0$ for all $j>k$; bandwidth makes this
automatic for $j>k+b$, so checking $j\in(k,k+b]$ suffices. Each pivot touches
an $\le b\times b$ window, giving $O(nb^2)$.
\end{proof}

\begin{remark}[The bordered extension]\label{rem:bordered}
$M(\lambda,t)$ has no bandwidth: its last row and column carry the generally
dense $\tfrac12(q_0-\sum_j\lambda_jq_j)$, so Lemma~\ref{lem:band} applied to
$M$ gives $O(n^3)$, not $O(Nd^3)$. The bordered sweep instead carries the
border as a \emph{separate} column, at one rank-one corner update per pivot.
That extension is \emph{asserted here, not proved}. Its soundness is pinned by
test in two ways: at the verdict level against the dense $M$-test, and pivot
by pivot against the dense $LDL^\top$ of $M$ at the shortest horizons of
\S\ref{sec:rpod}'s family. Its $O(Nd^3)$ count is read off the shipped sweep's
loop bounds --- each of the $n$ pivots touches a $\le b\times b$ band window
plus $\le b$ border-column entries and one corner update. No test measures the
bordered sweep's operation count; it is the band sweep's count that is
measured against its loop bound. The companion note records the extension's
status and catalogues the pinning receipts~\cite{leannote}.
\end{remark}

Both tests ship as trusted checkers in the Podium library. The
checkers work in exact rational arithmetic, refuse floats, and reject
malformed input --- non-rational entries, mismatched shapes, asymmetric
blocks --- before any verdict. One entry point is the band Hessian test. The
other is the corner-carrying bordered sweep: one band $LDL^\top$ carrying the
affine column, with the singular-$H$ range condition detected as a consistent
zero-pivot row, delivering the full $M(\lambda,t)\succeq0$ verdict in
$O(Nd^3)$ exact rational operations. Both entry points materialize the
$n\times n$ matrix before eliminating, so their memory --- and hence their
total work --- is quadratic in $n$; the $O(Nd^3)$ of Lemma~\ref{lem:band}
counts the sweep's rational \emph{operations}, which the assembly does not
change. A fixed-seed test suite pins both checkers against their dense
counterparts --- the band-$H$ sweep on the assembled $H$, the bordered sweep
on the assembled $M$, the latter being the dense $M$-test the library's
reference checker performs --- over randomized instances plus dedicated
singular-$H$ and zero-pivot cases. The companion note names the checkers, the
pinning suite, and the library commit that carries them~\cite{leannote}.

\begin{remark}[Coarse effectiveness]\label{rem:effective}
The sweep's pivots are ratios of consecutive leading minors of the bordered
matrix; with zero pivots skipped, the $k$-th pivot is a ratio of two minors
on the retained index set, so the identity survives the zero-pivot branch.
Clearing the common denominator that
$L_M=\mathrm{ht}\bigl(M(\lambda,t)\bigr)$ supplies, their bit-sizes are
\[
O\bigl((n{+}1)(L_M+\log(n{+}1))\bigr)
\]
by Hadamard's inequality and the classical exact-elimination argument of
Edmonds~\cite{edmonds1967}, and exact verification runs in time polynomial in
$n$ and $L_M$. The order $n{+}1$ is the worst case over the sweep, attained at
the \emph{bordered} $(n{+}1)\times(n{+}1)$ matrix itself: only the terminal
pivot's minors reach the border, the interior pivots' minors being leading
minors of the $H$-block alone, so bounding every pivot at the bordered order
and the bordered height is conservative. The implied constant is not $1$: a
pivot is a ratio of two minors, so after clearing, the instantiation of the
display is about twice the product. The height that governs here is the
bordered matrix's, not a per-stage height, and the difference is not
cosmetic. Assembly multiplies denominators across the coupling, so on the
demonstration family of \S\ref{sec:rpod} the bordered height dwarfs the
per-stage design height, and no per-stage height bounds even the
\emph{entries} of the matrix the sweep eliminates, let alone its pivots.
\S\ref{sec:rpod} measures those pivots inside the displayed expression formed
with $L_M$ at every horizon, and outside the one formed with the design
height $L$. Sharp growth laws are model-sensitive; \S\ref{sec:rpod} reports
measured growth.
\end{remark}

\noindent Remark~\ref{rem:effective}'s displayed expression takes both the
bordered height $L_M$ and the order $n{+}1$ of the bordered matrix because
what the sweep costs is bounded through the two together, and through neither
alone. Two admissible closing families witness this, each constructed by an
executable receipt catalogued in the companion note~\cite{leannote}. The
first family has bordered matrix entries all in $\{-1,0,1,2,3\}$,
$x^\star=0$, unit keep-outs contacted at $c_k=-1$, and multipliers
alternating $1$ and $0$. It holds $L_M=2$ at every horizon while its largest
interior pivot grows without bound --- $12$ bits at $n{+}1=9$ and $483$ at
$n{+}1=257$ --- so no function of $L_M$ alone bounds the cost. The second
family fixes the order at $n{+}1=9$ while its multipliers' denominators
grow: that raises $L_M$ from $10$ to $262$ bits and the largest interior
pivot from $85$ to $4116$ with it, so no function of the order alone bounds
the cost either.

The acceptance predicate --- a factorization with nonnegative pivots exists
iff the matrix is PSD --- is machine-checked in Lean in both directions; that
the shipped sweep computes such a factorization is not, and the companion note
records what is machine-checked and under which
hypotheses~\cite{leannote}.

\section{Conditional closure and the emitted certificate}\label{sec:theorem}
\begin{assumption}[Aggregate-PSD exactness]\label{a:psd}
At a KKT point $x^\star$ with multipliers $\lambda^\star\ge0$, the Lagrangian
curvature matrix $H(\lambda^\star)\succeq0$; global optimality of $x^\star$ is
then a \emph{conclusion} (Theorem~\ref{thm:main}), not a hypothesis, and A1 is
a solver-free PSD check on $H(\lambda^\star)$ that needs no prior knowledge of
optimality.
\end{assumption}
\begin{assumption}[Rational recoverable optimum]\label{a:rat}
$x^\star$ is rational: the active keep-out surfaces
$\lVert\Pi_jx_{\sigma(j)}-c_j\rVert=\rho_j$ meet the stationarity variety at a
rational point.
\end{assumption}

\noindent Assumption~\ref{a:psd} (A1 below; A2 is Assumption~\ref{a:rat}) is
the S-procedure exactness condition, and for the chain it is implied by the
solver-free weak-diagonal-dominance rule of \S\ref{sec:regime}.

\begin{remark}[Rational contact is designed, not generic]\label{rem:contact}
A2 is arithmetic-geometric and is arranged by construction rather than assumed
generically: when the active keep-out guards a sub-block of dimension greater
than one, the stationarity variety meets that active sphere at a generically
irrational point, so rational contact must be designed in. Designing it in
requires $\rho_j\in\Q$, since the named constructions --- Pythagorean
triples for $d=2$, stereographic projection of rational parameters for general
$d$ --- reach rational points only when $\rho_j$ is rational. At guard
dimension $2$ or $3$ a rational non-square $\rho_j^2$ can leave the sphere
with no rational point to design in at all: $x^2+y^2=3$ and
$x^2+y^2+z^2=7$ have none. Dimension $4$ or more always admits one, by
Lagrange's four-square theorem. A keep-out guarding a single coordinate
with rational radius $\rho_j\in\Q$ is the exception: contact there pins that
coordinate to the rational value $c_j\pm\rho_j$ directly. What
\eqref{eq:keepout}'s data fix is only $\rho_j^2\in\Q$: with $\rho_j^2$ a
rational non-square, single-coordinate contact is always irrational, while a
higher-dimensional sphere may still carry a rational point the named
constructions do not reach, as $(1,1)$ does on $x^2+y^2=2$.
Tolerance-free closure is thus a property of instances whose optimum is
rational; \S\ref{sec:limits} states what survives without A2.
\end{remark}

\begin{definition}[Constructed rational-sphere family]\label{def:constructed}
A \emph{constructed rational-sphere family} is a family of instances
of~\eqref{eq:qcqp} carrying a distinguished rational point $x^\star$, the
\emph{designed point} --- Theorem~\ref{thm:main}, not this definition, is what
makes it the optimum --- in which the following four conditions hold.
\begin{enumerate}
\item[(i)] Every keep-out is \emph{active} at $x^\star$, with radius
$\rho_j>0$, and each \emph{guarded} contact point $\Pi_jx^\star_{\sigma(j)}$
is placed at a rational point of its keep-out sphere via a Pythagorean or
stereographic parametrization, the remaining stage coordinates being placed
rationally as well, so that each stage's placed coordinates have height at
most $L$.
\item[(ii)] At every stage the active constraint gradients
$\nabla f_j(x^\star)$ are linearly independent. For one keep-out per stage
this is automatic, since $\rho_j>0$; for several keep-outs per stage it is a
condition on the contact geometry. Since those gradients are supported on the
one block, it caps the number of active keep-outs on a stage at $d$.
\item[(iii)] The multipliers $\lambda^\star_j\ge0$ are chosen rationally so
that, together with the stage's placed coordinates and keep-out centres, they
form a per-stage design tuple of height at most $L$, and so that
aggregate-PSD (A1) holds.
\item[(iv)] Every $W_k$ is invertible, and the goals $g_k$ are then recovered
from stationarity, so the $g_k$ --- and with them $q_0$ and $r_0$ --- are
\emph{derived} data, carried by the assembled height $L_{\mathrm{asm}}$
rather than by $L$.
\end{enumerate}
The definition forces no ordering between $L$ and $L_{\mathrm{asm}}$: on the
demonstration family (\S\ref{sec:rpod}) $L_{\mathrm{asm}}$ is by far the
larger, while an executable receipt exhibits an \emph{inverted} family,
admissible under this definition, on which it is the
smaller~\cite{leannote}. No coordinate is left inactive to
be solved for, so the emitted pair $(\lambda^\star,x^\star)$ is rational, of
per-stage height at most $L$ and hence of per-entry bit-size at most $2L$
(Definition~\ref{def:height}), by construction.
\end{definition}

\begin{theorem}\label{thm:main}\leavevmode
\begin{enumerate}
\item[(i)] \emph{Closure.} Under Assumptions~\ref{a:psd} and~\ref{a:rat},
with linear independence of the active constraint gradients (LICQ) at
$x^\star$, the exact-rational bracket closes: there exist rational
$(\lambda^\star\ge0,\,t^\star=J^\star,\,x^\star)$ with
$M(\lambda^\star,t^\star)\succeq0$, $f_j(x^\star)\ge0$ and
$t^\star=f_0(x^\star)=J^\star$ --- an accepting certificate in the sense of
Definition~\ref{def:protocol} --- with any number of keep-outs per stage, of
which LICQ leaves at most $d$ active on each.
\item[(ii)] \emph{Emitted size and recovery.} For a constructed
rational-sphere family (Definition~\ref{def:constructed}) the emitted
certificate $(\lambda^\star,x^\star)$ additionally has \emph{per-entry}
bit-size $O(L)$, with a constant depending on $d$ but not on $N$, and a
verifier holding the instance data and the transmitted witness $x^\star$
recovers the same $\lambda^\star$ within that bound, so the transmitted
$\lambda^\star$ need not be trusted.
\end{enumerate}
\end{theorem}
\begin{proof}
\emph{Closure.} By Assumption~\ref{a:psd}, $x^\star$ is a KKT point with
multipliers $\lambda^\star\ge0$; for a constructed family they are supplied
outright, by Definition~\ref{def:constructed}(iii)--(iv). A constraint
qualification is not what produces them, since the theorem does not assume
$x^\star$ is a local minimizer: optimality is the conclusion. By
A1 the Lagrangian $\mathcal L(x)=f_0(x)-\sum_j\lambda^\star_jf_j(x)$ is
convex and $x^\star$ is a stationary point, hence a global minimiser:
$\min_x\mathcal L=\mathcal L(x^\star)=f_0(x^\star)$, the last step by
complementary slackness. Thus $M(\lambda^\star,t^\star)\succeq0$ at
$t^\star=f_0(x^\star)$, and weak duality gives $t^\star\le J^\star$, while
feasibility gives $J^\star\le f_0(x^\star)=t^\star$. Hence
$t^\star=J^\star=f_0(x^\star)$.\footnote{The mechanized counterparts of the
weak-duality step and the attainment leg, and the hypotheses those theorems
take, are itemized in the companion note~\cite{leannote}.} Nothing in this
argument uses one keep-out per stage: with several active keep-outs on a
stage the multipliers enter $\mathcal L$ additively and complementary
slackness holds term by term. Two exact receipts with two active keep-outs
on one stage accompany this~\cite{leannote}: one with disjoint-coordinate
contacts, closing at $t^\star=169/320$, and one whose contacts' gradients
share their whole support, closing at $t^\star=5773/3600$.

\emph{Rationality of the multipliers.} The data and $x^\star$ are rational,
so stationarity $\nabla f_0(x^\star)=\sum_j\lambda_j\nabla f_j(x^\star)$
together with complementary slackness ($\lambda_j=0$ at every inactive $j$)
is a rational linear system in the active multipliers, and under LICQ ---
item (ii) of Definition~\ref{def:constructed} for the constructed families
--- it has a unique solution, which is therefore rational.

\emph{Bit-size of the emitted pair.} A2 alone does not bound the height: for
quadratic optimization at large the KKT system is polynomial, and
Bienstock-type gadgets satisfy A2 with $2^n$-bit
optima~\cite{bienstock2020}. In a constructed family the design order of
Definition~\ref{def:constructed} settles the emitted pair directly:
$x^\star$ by (i) and $\lambda^\star$ by (iii) are \emph{placed} rationally at
per-stage height at most $L$, and only then are the goals recovered by (iv).
The placement leaves nothing to solve for: every keep-out being active, the
guarded coordinates sit on their keep-out sphere and the rest are placed
freely, with no inactive coordinate left over. So
$(\lambda^\star,x^\star)$ has per-entry bit-size $O(L)$, independent of $N$.

\emph{What the verifier reconstructs.} What needs an argument is that
$\lambda^\star$ is exactly what a verifier holding the instance data and the
transmitted witness $x^\star$ reconstructs. Fix a stage $s$, let $A_s$ be the
keep-outs active there and $m_s=|A_s|$; their gradients
$\nabla f_j(x^\star)=2\Pi_j^\top(\Pi_jx^\star_s-c_j)$ are supported on block
$s$, hence live in $\Q^d$, and have Euclidean norm $2\rho_j>0$, and by (ii)
they are linearly independent, so $m_s\le d$. Write $G_s\in\Q^{d\times m_s}$
for their block-$s$ restrictions and $v_s=[\nabla f_0(x^\star)]_s$.
Stationarity reads $G_s\lambda^\star_{A_s}=v_s$ with $G_s$ of full column
rank, so the normal equations
\[
\bigl(G_s^\top G_s\bigr)\lambda^\star_{A_s}=G_s^\top v_s
\]
have $\lambda^\star_{A_s}$ as their unique solution, and the verifier
reconstructs the emitted multipliers rather than trusting them. What it reconstructs is the
\emph{placed} $\lambda^\star$ itself, so
Definition~\ref{def:constructed}(iii) together with
$\mathrm{bit}\le2\,\mathrm{ht}$ (Definition~\ref{def:height}) bounds it by
$\mathrm{bit}(\lambda^\star_j)\le2L$ --- the design order, not the solve, is
what makes the recovered object small. Nothing here asks the active gradients
on a stage to have disjoint supports: the shared-support receipt above
exhibits a Definition-\ref{def:constructed}-admissible stage on which they
share every coordinate, and checks that the normal-equation solve
recovers $\lambda^\star$ there. What that bound does \emph{not} cover is the
cost of the same solve run from the assembled instance alone. That solve's
\emph{operands} stay at the design height on every family measured, and not
by accident: stationarity makes the objective gradient
$[\nabla f_0(x^\star)]_s$ the multiplier combination of the active gradients
$2\Pi_j^\top(\Pi_jx^\star_s-c_j)$, which are design-tuple data. What carries
the assembled height is the \emph{formation}: a verifier holding only the
assembled coefficients and the transmitted witness builds
$[\nabla f_0(x^\star)]_s$ as $2(P_0x^\star)_s+q_0^{(s)}$. That formation is
bounded through the assembled height together with the height of the placed
coordinates and active centres, not through $L$ alone, and which of the two
dominates it is a property of the instance; \S\ref{sec:rpod} reports the
measured figures. Acceptance recomputes $t^\star$ and runs the exact
$LDL^\top$; the size of those intermediates is a \emph{verification} cost
(Remark~\ref{rem:effective}), not a property of the emitted pair.
\end{proof}

\begin{remark}[Size of the certified value]\label{rem:value}
Write $\den(\cdot)$ for the lcm of the denominators of a tuple of rationals.
At a \emph{closing} certificate --- where complementary slackness makes the
Lagrangian value the verifier recomputes equal to $f_0(x^\star)$ --- the
certified value $t^\star=f_0(x^\star)$ is a sum of $N$ stage terms, the
$k$-th the quadratic form in $x^\star$ carried by the \emph{assembled} stage
coefficients $(P_0^{(k)},q_0^{(k)},r_0^{(k)})$: block row $k$ of $P_0$,
block $k$ of $q_0$, and a share $r_0^{(k)}$ of the constant,
block-tridiagonality making that row reach only stages $k{-}1$, $k$ and
$k{+}1$. A merely bracketing $t$ additionally carries $\den(\lambda)$ and the
keep-out coefficients' $\den(q_j)$ and $\den(r_j)$, with
$q_j=-2\Pi_j^\top c_j$ and $r_j=\lVert c_j\rVert^2-\rho_j^2$ read
off~\eqref{eq:keepout}, so that $\den(r_j)=\den(\rho_j^2)$ for the
origin-centred keep-outs of \S\ref{sec:rpod} and not in general. Any
splitting of the constant with $\sum_kr_0^{(k)}=r_0$ and
$\den(r_0^{(k)})\mid\den(r_0)$ serves --- for instance carrying all of $r_0$
on one stage --- and the receipt takes the conservative $\den(r_0)$ at every
stage. In four steps:
\begin{enumerate}
\item[(1)] Each summand's denominator divides
$\den(P_0^{(k)},q_0^{(k)},r_0^{(k)})\cdot\den(x^\star)^2$, the form being
quadratic in the witness, and a finite sum's denominator divides the lcm of
the summands', so
\[
\den(t^\star)\ \Big|\ \lcm_{k\le N}
\den\bigl(P_0^{(k)},q_0^{(k)},r_0^{(k)}\bigr)\cdot\den(x^\star)^2;
\]
hence $\den(t^\star)$ is bounded in $N$ whenever those denominators come
from a bounded set, as all families of Table~\ref{tab:results} do.
\item[(2)] Stated against the heights of \S\ref{sec:sizes} --- the
\emph{coefficient} height $L_{\mathrm{asm}}$, not the bordered $L_M$, since
it is the coefficient denominators that appear above --- the right-hand side
divides the common denominator that $L_{\mathrm{asm}}$ supplies times
$\den(x^\star)^2$, so
$\mathrm{bitlen}(\den(t^\star))\le L_{\mathrm{asm}}+2\,\mathrm{ht}(x^\star)$.
Here $\mathrm{ht}(x^\star)$ is taken over all stages at once, whereas $L$
bounds a stage's design tuple one stage at a time, so the more compact
$\mathrm{bitlen}(\den(t^\star))\le L_{\mathrm{asm}}+2L$ additionally asks
that the placed coordinates share one denominator along the horizon --- as
they do on the family of \S\ref{sec:rpod}.
\item[(3)] The numerator adds only the $\log N$ of the sum over stages and
the log of the number of monomials one block row contributes --- at most
$3d^2{+}d{+}1$, since block-tridiagonality leaves block row $k$ at most $3d$
nonzero columns across its $d$ rows, plus $q_0$'s block of $d$ and the
constant --- giving
$\mathrm{bit}(t^\star)\le2(L_{\mathrm{asm}}+2L)+\lceil\log_2N\rceil
+\lceil\log_2(3d^2{+}d{+}1)\rceil=O(L_{\mathrm{asm}}+L+\log N)$, so whenever
$L_{\mathrm{asm}}$ is bounded in $N$, as it is on the demonstration family
of \S\ref{sec:rpod}, the value's size grows at most logarithmically along
the horizon.
\item[(4)] The governing quantity is \emph{not} the lcm of the raw data
denominators, and cannot be: the assembled coefficients are products and
quotients of the data --- $A_k^\top R_kA_k$ multiplies two denominators, and
recovering $g_k$ from stationarity (Definition~\ref{def:constructed}(iv))
divides by $W_k$, so a data \emph{numerator} can become a coefficient
denominator --- after which the goal term squares them. On the CW family of
\S\ref{sec:rpod} the raw-data lcm is $130$ bits while $\den(t^\star)$ is
$450$--$526$ bits; the displayed divisibility holds at every one of the $23$
horizons $N=2,\dots,24$, and $\den(t^\star)$ divides the raw-data lcm at
none of them.
\end{enumerate}
\S\ref{sec:rpod} checks all three inequalities, and the divisibility step
behind them, at every horizon of the demonstration family.
\end{remark}

\noindent Why Remark~\ref{rem:value} is stated against the heights it names,
and why its bit bound carries a horizon term, is settled by admissible
closing families, each constructed by an executable receipt~\cite{leannote}.
The certified value's \emph{denominator} is bounded through
$L_{\mathrm{asm}}$ together with the height of the placed coordinates, but
its \emph{bit-size} additionally needs the horizon: an admissible closing
family holds $L_{\mathrm{asm}}=2$ and $\mathrm{ht}(x^\star)=1$ at every
horizon while $t^\star=N$ exactly. Nor is either quantity bounded through
$L$ alone: integer goal weights of growing magnitude take $\den(t^\star)$
unbounded at fixed $L=1$, since goal recovery divides by $W_k$ and the
weight's numerator lands in the value's denominator.

No unconditional polynomial bit bound exists: quadratic optimization
\emph{at large} can force $2^n$-bit exact solutions~\cite{bienstock2020},
and small-rational-solution existence is strongly NP-hard. What we
\emph{can} support is Theorem~\ref{thm:main}'s guarantee for the constructed
families, and the reason it is supportable is the design order of
Definition~\ref{def:constructed}: $(\lambda^\star,x^\star)$ is placed at
per-stage height $O(L)$ before any datum is derived from it, so no solve
stands between the design tuple and the emitted pair.

\section{Diagonal dominance as a solver-free benign criterion}\label{sec:regime}
We call an instance \emph{benign} when aggregate-PSD (A1) holds at its optimum
--- the S-procedure is then exact and the bracket can close. This is a
property of the instance, not of a solver; this section gives a sufficient,
solver-free condition checkable stage-by-stage along the sweep.

\begin{lemma}[Solver-free benign criterion]\label{lem:dd}
Call $H$ \emph{weakly diagonally dominant} (WDD) when
$H_{ii}\ge\sum_{j\ne i}|H_{ij}|$ on the \emph{scalar} entries of the
assembled $H$, not its $d\times d$ blocks; the left side is $H_{ii}$, not
$|H_{ii}|$, so the inequality already forces a nonnegative diagonal. If $H$
is symmetric and WDD, then $H\succeq0$, and the Schur complement after each
scalar pivot of the sweep is again WDD, so every pivot along the sweep is
nonnegative.

For the identity-coupled chain ($A_k=I$, $R_k=I$) with full-stage keep-outs
($\Pi_j=I$), WDD is exactly the per-stage condition $w_k\ge\lambda_k$, where
$w_k$ is the scalar goal weight ($W_k=w_kI$) and $\lambda_k$ the total
multiplier load on stage $k$: the goal weight dominates the keep-out
multipliers. Under a projection $\Pi_j$ the load falls only on the guarded
coordinates, so WDD becomes the per-coordinate condition
$w_k\ge\max_i\lambda_{k,i}$, with $\lambda_{k,i}$ the summed multipliers of
the keep-outs guarding coordinate $i$ of stage $k$. The stage condition
$w_k\ge\lambda_k$ remains sufficient, and remains sharp whenever some
coordinate is guarded by every keep-out on the stage --- in particular for a
single projected keep-out; it is strictly conservative exactly when the
guarded sets of the keep-outs carrying \emph{positive} multipliers share no
coordinate, which, those sets being nonempty, forces at least two positive
multipliers.
\end{lemma}
\begin{proof}
A symmetric WDD matrix is PSD by Gershgorin (machine-checked;
\cite{leannote}). The Schur complement of a strictly diagonally dominant
matrix is strictly diagonally dominant~\cite{carlson1979}; the weak case
follows by perturbing the diagonal by $\varepsilon\to0^+$, zero pivots having
zero rows in the symmetric WDD case and being skipped, as in
Lemma~\ref{lem:band}. Applying this one pivot at a time keeps every Schur
complement WDD, hence every pivot nonnegative; $H\succeq0$ itself is the
Gershgorin step above. In the identity-coupled scalar specialization the whole
check reduces to the local test $w_k\ge\lambda_k$ on each row of $H$ ---
interior rows $[-1,\,2+w_k-\lambda_k,\,-1]$ and the two boundary rows
$[\,1+w_k-\lambda_k,\,-1\,]$ --- with no block pivot $S_k$ computed, which
also avoids presuming the nonsingularity that $S_k$'s recursion would need.
\end{proof}

WDD is sufficient, not necessary. The one case where it is also necessary is
special: for the identity-coupled free-endpoint chain with \emph{uniform}
excess $w_k-\lambda_k\equiv c$, the path-Laplacian critical mode is the
constant vector, so $H\succeq0$ iff $c\ge0$. That is a boundary statement,
and its hypotheses --- identity coupling, uniform excess --- are
load-bearing.

Nor does one keep-out per stage guarantee A1:
\begin{example}[Certified duality gap]\label{ex:gap}
A $d{=}1$, $N{=}2$ instance of~\eqref{eq:qcqp} with one
keep-out per stage, stage-coupling coefficient $-3/2$, and nonzero goals has
$J^\star=212/13$, by exact enumeration of the branch-KKT candidates. A
rational moment point caps \emph{every} level-one certificate at $815/64$ by
weak duality. The level-one gap is thus certified $\ge2973/832\approx3.57$,
entirely over $\Q$.
\end{example}
\noindent The same receipt discharges the step from a gap to a failure of A1
rather than leaving it to the reader: $x^\star$ is rational (A2) and exactly
one keep-out is active there, so LICQ is automatic and Theorem~\ref{thm:main}
would close the bracket if A1 held. It does not: stationarity recovers
$\lambda^\star=24/13$ and $H(\lambda^\star)$ has determinant $-7/4$, so
$H(\lambda^\star)\not\succeq0$. This is why aggregate-PSD (A1) is a real
assumption and not automatic: the block-banded smoothness coupling is not
itself a gap source under WDD, but general coupling with nonzero goals can
break it.

\paragraph{Hard dynamics.} Hard equality dynamics $x_{k+1}=A_kx_k+B_ku_k$ are
outside~\eqref{eq:qcqp} by definition: they add control variables to a
decision vector \S\ref{sec:setup} takes to be the stage states alone, and an
equality to a class that carries only inequalities.
Lemma~\ref{lem:dd}'s criterion is likewise stated for the penalty-coupled
assembly only. The \emph{closure} argument is less fragile than that: an
affine equality contributes nothing to the Lagrangian curvature matrix and
enters weak duality with a free multiplier, so wherever the analogue of
aggregate-PSD holds at a rational KKT point the same constant-multiplier
bracket closes. An exact receipt certifies both halves over
$\Q$ on $d{=}1$ instances with one keep-out and one
equality~\cite{leannote}. On the first, the best constant-multiplier bound is
exactly $1$ while $J^\star=3/2$, a certified gap of $1/2$. The diagnosis,
though, is Example~\ref{ex:gap}'s, not a statement about hard coupling:
stationarity there recovers $\lambda^\star=3/2$ against a unit curvature, so
$H(\lambda^\star)=\mathrm{diag}(-1/2,1,1)$ is not PSD and A1 fails, with A2
and LICQ intact. On its sibling --- the same variables and the same equality,
with the goal moved and the keep-out radius doubled --- stationarity recovers
$\lambda^\star=3/4$ and $H(\lambda^\star)=\mathrm{diag}(1/4,1,1)\succeq0$, so
A1 holds and the constant-multiplier bracket closes exactly at
$t^\star=J^\star=9/4$: the same trusted dense $M$-test accepts $t^\star$ and
refuses $t^\star+10^{-6}$. The exhibited gap therefore measures A1, as
Example~\ref{ex:gap}'s does, and no conclusion about Theorem~\ref{thm:main}
follows from it.

\section{Demonstration: an inverse-designed CW-coupled instance}\label{sec:rpod}
We instantiate~\eqref{eq:qcqp} as a \emph{certification stress instance with
Clohessy--Wiltshire structure} --- not a flyable maneuver. Consecutive stages
are coupled through the planar CW state-transition matrix $\Phi$ (mean motion
$n_0=10^{-3}$\,rad/s and sampling phase $\tau=n_0\,\Delta t=1/2$, so
$\Delta t=500$\,s) via the soft-dynamics \emph{full-state penalty}
$\lVert x_{k+1}-\Phi x_k\rVert^2_R$:
there is no control variable and no thrust model, the keep-out radius cycles
$\{1/2,1,3/2\}$ per stage, and the instance is \emph{constructed by inverse
design} (each
stage's guarded contact placed on a rational keep-out-sphere point, goals
recovered from stationarity, benign multipliers chosen so A1 holds). The stage
state is $x_k=(x,y,v_x,v_y)\in\Q^4$ and the keep-out guards its \emph{position}
sub-block: $\Pi_k$ in~\eqref{eq:keepout} selects $(x,y)$, giving
$\lVert\Pi_kx_k\rVert^2-\rho_k^2\ge0$ about the target at the LVLH origin, with
the two velocity coordinates placed rationally by the design. It demonstrates
the certification pipeline on a rational-optimum instance whose \emph{coupling
matrix} is a genuine CW state-transition matrix; the \emph{trajectory} is not
CW-consistent and is not meant to be.

A stage is one sampled state: the horizon variable $N$ counts state samples
taken $\Delta t=500$\,s apart, and one sampling step advances the reference
orbit's phase by $\tau=1/2$\,rad, so about $12.6$ steps span one orbit. The
$N$ stage states are joined by $N-1$ coupling residuals, so a horizon spans
$(N-1)\,\Delta t$. At the largest demonstrated horizon, $N=24$, that is
$11500$\,s: about $3.2$\,h, or $1.8$ orbits in the model's units. The same
horizon carries $n=Nd=96$ scalar variables at $d=4$ per stage. The translation to
time is specific to the CW-coupled family: the shared-state families of
Table~\ref{tab:results} are abstract chains with no sampling interval.

The design places the optimum where stationarity requires it: at $N{=}24$ the
per-step dynamics residual $\lVert x^\star_{k+1}-\Phi x^\star_k\rVert$ runs
$134.8$--$244.1$ against states of norm $0.559$--$1.572$, and the goals
recovered from stationarity reach $\lVert g_k\rVert\approx1.01\times10^{5}$.
Units are nominal --- positions in metres,
velocities in metres per second --- so $R=I_4$ and $W_k=w_kI_4$ are not
dimensionally homogeneous weights and the norms just quoted mix the two,
consistent with the instance not being a flyable maneuver. So this does not
discover an optimum for an operational CW problem, and the goal term is a
goal-attraction penalty, not a tracking objective. The CW state-transition
matrix is rationalized entrywise to denominators at most $10^5$, giving
per-entry bit-size at most $42$. The rationalization breaks the symplectic
structure: after rounding, $\det\Phi=1+1.2\times10^{-8}$ exactly, against
$\det\Phi=1$ for the CW flow. The certificate is therefore a theorem about
the \emph{rational model only}: it certifies optimality and keep-out
satisfaction for the rationalized dynamics, not for the continuous CW flow.
The rounding moves each entry by at most $2.8\times10^{-10}$ in absolute
value and $2.0\times10^{-7}$ relative; it is the $10^5$ denominator cap, not
the $42$-bit entry size, that sets those figures.

Everything is checked in $\Q$ by the three sweeps of \S\ref{sec:riccati} ---
the band-$H$ test, the corner-carrying bordered sweep on the full $M$, and the
dense $M$-test --- which agree on every row of Table~\ref{tab:results}. At
every one of the $23$ horizons $N=2,\dots,24$ the bracket closes, so weak
duality already forces $g(\lambda^\star)=t^\star$ and a terminal corner pivot
of exactly $0$; as a sanity check the bordered sweep refuses $t^\star+10^{-6}$
at every horizon, and the guarded position sub-blocks of the optima lie
exactly on the keep-out spheres.

\begin{table}[t]\centering
\caption{Exact-rational trajectory certificates (all verdicts exact over $\Q$).
``emitted bits'' is the worst-entry bit-size of the emitted
$(\lambda^\star,x^\star)$; ``value bits'' is the size of the certified value
$t^\star$ the verifier recomputes (Remark~\ref{rem:value}). The size column
carries both units: $N$ is the number of stages, $d$ the per-stage state
dimension, and $n=Nd$ the number of scalar variables. The regime column
places each family in the vocabulary of \S\ref{sec:regime}: \emph{benign}
means A1 holds at the optimum, \emph{A1 violated} means it fails and the
checker refuses. Rows 2 and 3 are constructed rational-sphere families
(Definition~\ref{def:constructed}); row 1 is not --- the scalar shared-state
family leaves inactive stages with solved coordinates. It is benign, and
covered by Theorem~\ref{thm:main}'s closure leg, but not by its $O(L)$
bit-size bound, which is stated only for constructed families; so its two bit
columns are empirical measurements. The CW row's ranges are over all $23$
horizons $N=2,\dots,24$; rows 1 and 2 are over the horizons the measurement
receipt samples, up to the max size shown.}
\label{tab:results}
% \small + tighter \tabcolsep: the dual-unit size column pushed the natural
% width past the measure (38pt overfull); check_refs fails on any overfull
% box, so the table must fit.
\small
\setlength{\tabcolsep}{4pt}
\begin{tabular}{@{}llcccc@{}}
\toprule
Family & regime & max size & closes & emitted bits & value bits\\
\midrule
Scalar shared-state & benign & $N{=}129$ ($n{=}129$) & yes & $13$--$18$ & $13$--$26$\\
Vector shared-state ($d=1..4$) & benign & $n{=}Nd{=}96$ & yes & $3$--$4$ & $15$--$32$\\
CW-coupled keep-out & benign & $N{=}24$ ($n{=}96$) & yes & $6$--$8$ & $931$--$1089$\\
Aggregate-PSD violated & A1 violated & --- & no (refuses) & --- & ---\\
\bottomrule
\end{tabular}
\end{table}

The emitted pair is small in every family --- $18$ bits or fewer --- while the
recomputed CW \emph{value} ($\sim\!10^3$ bits) comes from the inverse design,
not from the raw data. The goal term $r_0=\sum_kw_k\lVert g_k\rVert^2$
\emph{squares} the denominators recovered from stationarity: $\den(t^\star)$
and $\den(r_0)$ coincide at every horizon $N\ge6$ (both $519$ bits at
$N{=}24$), and $\den(t^\star)$ divides $\den(r_0)$ at $22$ of the $23$
horizons --- the exception is $N{=}5$, where the two differ by a factor $6$
the other way. Per Remark~\ref{rem:value} it is the \emph{denominator} that
is bounded in $N$: $\den(t^\star)$ stays inside $450$--$526$ bits at every
horizon out to $N=6143$. The bit-size keeps the remark's residual $\log N$
drift, measured at $931$--$1089$ bits over all $23$ demonstration horizons.
The verifier's multiplier reconstruction
(Theorem~\ref{thm:main}) shows the same split between the design and assembled
heights. The solve's \emph{operands} stay at the design height: $5$--$6$
bits against $L=6$ on this family, and $259$ bits against $L=258$ on the
inverted family noted after Definition~\ref{def:constructed}. They stay
there because stationarity makes the objective gradient the multiplier
combination of the active keep-out gradients, which are design-tuple data.
The \emph{formation} of that gradient
from the assembled coefficients carries the assembled height instead: on this
family its two terms run $242$--$276$ bits against $L=6$ and cancel to the
design height, while on the inverted family they do not cancel, the formed
gradient staying at $259$ bits against $L_{\mathrm{asm}}=3$--$8$.

\paragraph{The heights, measured.} The four separate heights of
\S\ref{sec:sizes} are ``in general different numbers'' on this family in the
strongest sense: at every one of the $23$ horizons its four tuples carry four
pairwise distinct heights, in the strict order
$\mathrm{ht}(x^\star)<L<L_M<L_{\mathrm{asm}}$; Table~\ref{tab:measured}
reports the measured ranges.

\begin{table}[ht]\centering
\caption{The four heights of \S\ref{sec:sizes}, measured on the CW-coupled
demonstration family over all $23$ horizons $N=2,\dots,24$; they are pairwise
distinct at every horizon, in the strict order
$\mathrm{ht}(x^\star)<L<L_M<L_{\mathrm{asm}}$.}
\label{tab:measured}
\begin{tabular*}{6.1in}{@{}p{1.15in}p{2.35in}p{2.2in}@{}}
\toprule
Height & measured & what it governs here\\
\midrule
$\mathrm{ht}(x^\star)$ & $5$ bits at every horizon &
Remark~\ref{rem:value}'s sharper first inequality\\
\addlinespace
$L$ & $6$ bits, attained at every horizon (per-stage values $3$, $5$ or $6$)
& Theorem~\ref{thm:main}'s $O(L)$ leg; Table~\ref{tab:results}'s emitted
$6$--$8$ bits\\
\addlinespace
$L_M$ & $277$--$284$ bits & Remark~\ref{rem:effective}'s pivot expression\\
\addlinespace
$L_{\mathrm{asm}}$ & $483$--$563$ bits & Remark~\ref{rem:value}'s bounds on
the certified value; carries the $519$-bit $\den(r_0)$\\
\bottomrule
\end{tabular*}
\end{table}

The per-stage \emph{design} tuple of Definition~\ref{def:constructed}(iii) ---
a stage's placed coordinates, its keep-out centre (the LVLH origin here) and
its multiplier --- has height at most $L=6$ bits at every stage of every
horizon, the maximum attained at every one of them. Table~\ref{tab:results}'s
emitted $6$--$8$ bits sit next to that figure by construction rather than by
measurement: the emitted pair is part of the very tuple $L$ is taken over, so
$\mathrm{bit}(\cdot)\le2L$ (Definition~\ref{def:height}) holds identically,
here and on any constructed family. What the horizons do show is the point
actually at issue: neither figure grows with $N$. The assembled
coefficients $(P_0,q_0,r_0,g)$ have height $L_{\mathrm{asm}}=483$--$563$ bits,
consistent with Definition~\ref{def:constructed}(iv), which makes $r_0$
derived data of the assembled height rather than of $L$: the $519$-bit
$\den(r_0)$ above is carried by $L_{\mathrm{asm}}$ and could be carried by
neither the design height nor $L_M$. And the assembled bordered matrix
$M(\lambda^\star,t^\star)$ that the sweep eliminates has height
$L_M=277$--$284$ bits, the smaller of the two assembled figures because its
corner $r_0-\sum_j\lambda^\star_jr_j-t^\star$ cancels those two large
denominators against each other.

Instantiating Remark~\ref{rem:effective}'s expression at the bordered order
$n{+}1$ gives $(n{+}1)\bigl(L_M+\lceil\log_2(n{+}1)\rceil\bigr)$, running
$2529$--$28227$ bits over the horizons. The measured pivots run from $777$
bits at $N{=}2$ to $17893$ at $N{=}24$ and stay inside that expression at
all $23$ horizons. Because the
bracket closes, the terminal corner pivot $g(\lambda^\star)-t^\star$ is
exactly $0$ at every horizon, so the figure reported is the largest
\emph{interior} pivot of the corner-carrying sweep. The same expression formed
with the design height $L$ runs only $90$--$1261$ bits and is exceeded at all
$23$ horizons. The pair of measurements $L_M=277$--$284$ bits against $L=6$
already says that on this family no per-stage height bounds even the
\emph{entries} of the bordered matrix, let alone its pivots. Here it is the
bordered height, and not a per-stage height, that
Remark~\ref{rem:effective}'s expression must be instantiated with.

That ordering is a property of the CW assembly and not of
Definition~\ref{def:constructed}, which forces neither direction: an exact
receipt exhibits an admissible family, closing on all three trusted
checkers, on which $L_M\le8$ bits stands against a design height
$L=257$--$258$, so there the per-stage design height bounds every entry of
the bordered matrix~\cite{leannote}. Remark~\ref{rem:value}'s three inequalities hold at all
$23$ horizons as well --- the compact
$\mathrm{bitlen}(\den(t^\star))\le L_{\mathrm{asm}}+2L$, the value bound
$\mathrm{bit}(t^\star)\le2(L_{\mathrm{asm}}+2L)+\lceil\log_2N\rceil
+\lceil\log_2(3d^2{+}d{+}1)\rceil$, and the sharper first form
$\mathrm{bitlen}(\den(t^\star))\le L_{\mathrm{asm}}+2\,\mathrm{ht}(x^\star)$
--- as does the divisibility step behind them. The pivots themselves grow
measurably linearly in the horizon --- $780$ bits at $N{=}2$ to $17889$ at
$N{=}24$, measured over the blocks $S_k$ of~\eqref{eq:riccati} rather than
over the scalar pivots, hence a few bits either side of the bordered figures
above; this is reported as data.

\paragraph{Reproducibility.} Each row of Table~\ref{tab:results} is
constructed by a named driver in exact rational arithmetic --- the three
closing rows certified, the refusal row refused --- and each driver
verifies its emitted certificate through the exact dense $M$-test. The
refusal row's instance takes $\lambda_k>w_k$ at $N{=}4$, $d{=}2$.
Measurement receipts recompute both bit columns at the sizes
Table~\ref{tab:results} headlines, re-verify every closing row through all
three trusted routes (dense $M$-test, band-$H$ sweep, bordered-$M$ sweep),
and recompute every denominator statistic quoted in this section at all $23$
horizons; a shipped test suite executes each receipt and asserts its
published constants. The trusted checkers refuse floats, and the CW
state-transition matrix alone is synthesized in floating point and
rationalized before any certification, as above. The receipts accompany this
record as supplementary material, and the trusted checkers ship in the
Podium repository~\cite{oltean2026podiumsw}, described
in~\cite{oltean2026podium}. The companion note catalogues every receipt by
role and records the bundle's conventions, including the bit-size convention
and the library commit that carries the band checkers~\cite{leannote}.

\section{Related work and positioning}\label{sec:related}
\paragraph{Exact-rational certificates for optimization.} The direct
antecedent of tolerance-free certificates for optimization outputs is the
exact rational LP/MIP line: VIPR's machine-verifiable exact certificates for
integer-programming results~\cite{cheung2017vipr}, exact rational MIP solving
in SCIP~\cite{eifler2021exact}, and the exact LP solver
QSopt\_ex~\cite{applegate2007qsopt}; the bit-size control of exact
elimination goes back to Edmonds~\cite{edmonds1967}. Certificates measured by
their verification complexity are the program of~\cite{dumas2014}. Ours
extends the exact-certificate discipline from linear to a structured
nonconvex quadratic class, with the banded-KKT/Riccati structure doing the
work.

\paragraph{Trajectory SDP relaxations and tightness.} Sparse moment--SOS for
trajectory optimization~\cite{strom2024} and
perception~\cite{yang2023perception} is the closest numerical program:
floating-point, tolerance-based certificates --- STROM reports certified
suboptimality below $1\%$ for most of its five trajectory problems, near
$1\%$ for the hardest (car back-in), and Yang and Carlone report a sparse
relaxation that is \emph{empirically exact}, deeming a rounded solution
globally optimal at a relative suboptimality below $10^{-3}$ --- obtained
from level-two (sparse second-order) relaxations that close gaps left open at
level one, with tightness observed rather than proved. The two lines are
complementary: exact-arithmetic acceptance on benign, rational-optimum
(constructed) instances versus broad numerical coverage. On the same
constraint geometry as \S\ref{sec:regime}, Graesdal et
al.~\cite{graesdal2026} give necessary \emph{and} sufficient tightness
conditions for a \emph{single}-sphere-obstacle trajectory SDP relaxation;
tightness in their own many-obstacle experiments is measured empirically, not
obtained from that theorem. Ours is only sufficient and holds with any number
of keep-outs per stage --- of which LICQ leaves at most $d$ active on each
--- but is solver-free and checkable in exact arithmetic along the sweep, and
their single-obstacle conditions delimit what a sharp benign boundary could
look like. In the adjacent domain of spacecraft conjunction avoidance, Vega
et al.~\cite{vega2025} solve collision-avoidance exclusion QCQPs via a Shor
relaxation observed empirically tight --- on a different formulation, with
control variables and conjunction-derived constraints, so not a measurement
of A1 on our class. We share no benchmark with any of these, so this
paragraph is positioning, not comparison. A general tightness framework for
QCQP SDP relaxations~\cite{wang2021tightness}, graph-structural exactness
conditions~\cite{sojoudi2014}, exactness on forest
structures~\cite{azuma2020}, and local-to-global exactness from clique-wise
sub-SDPs~\cite{kojima2026} position A1 within a broader theory (a path is a
tree, but our concave keep-outs need not satisfy their hypotheses); chordal
conversion and PSD completion~\cite{zhang2023,grone1984} address the
structural SDP side. Lossless convexification~\cite{acikmese2011} exactly
convexifies nonconvex \emph{control}-set constraints via the maximum
principle --- complementary in mechanism to exact certificates for nonconvex
\emph{state} constraints.

\paragraph{Rational recovery and exact SDP.} Exact rational SOS/PSD
certificate recovery is due to~\cite{peyrl2008,kaltofen2012}, with
RealCertify the standard tool~\cite{magron2018realcertify}: its multivariate
route computes an approximate decomposition of a perturbed input with an
\emph{arbitrary-precision} SDP solver (SDPA-GMP) and recovers an exact
rational one symbolically, while its univariate \texttt{univsos1} uses no SDP
at all --- a recursive construction over $\Q$ from root isolation and
quadratic under-approximations. Exact algorithms for LMIs and the SPECTRA
library are the exact-SDP lineage~\cite{henrion2019spectra}. Bit-size bounds
for exact SOS certificates --- univariate~\cite{magron2018},
multivariate~\cite{magron2021multivariate,davispapp2023} --- do not exploit
band structure, an orthogonal axis to Theorem~\ref{thm:main};
machine-checkable SOS certification is developed in~\cite{magron2015formal},
the Coq-interfaced companion to the template method of~\cite{allamigeon2014}.

\paragraph{Verified solvers and a-posteriori numerics.} A complementary line
proves the \emph{solver} correct before it runs (credible autocoding and
verified convex-optimization codes~\cite{wang2014,davy2018,cohen2020}); Roux
et al.~\cite{roux2018} validate solver output a posteriori via a
floating-point Cholesky with proved rounding-error bounds --- rigorous but
tolerance-based. The a-priori and a-posteriori guarantees compose with ours
rather than compete. On \emph{why exactness}: a rigorous inequality
certificate need not be exact. Interval arithmetic or a verified
Cholesky~\cite{jansson2009,roux2018} certifies $H(\lambda)\succ0$ with a
positive margin, hence a rigorous lower bound, at fixed precision. A
positive-margin criterion cannot reach the semidefinite boundary, and the
closing certificates of Theorem~\ref{thm:main} sit exactly there: at closure
the corner Schur complement is exactly $0$, and singular $H$ is the case
\S\ref{sec:setup} handles through the consistent zero-pivot row. Fixed
precision reaches the boundary only in a narrow case: the operations
witnessing the boundary contact must be exact in the working precision, so
that the boundary quantity is enclosed with zero width. Pivots away from the
boundary may round, their positive-width enclosures absorbed by their
positive margins. Even in that case a Cholesky does not certify the
boundary, because it stops at an exactly computed zero pivot: on the
singular $H$ with rows $(1,-1)$ and $(-1,1)$, where an $LDL^\top$ and the
symmetric eigensolver are exact, LAPACK's \texttt{dpotrf} declines the
factorization, reporting the order-2 leading minor not positive definite.
The rationalized CW instance of \S\ref{sec:rpod} is not of that
boundary-case kind at all: its entry denominators are not powers of two.
Exact receipts pin each of these boundary facts; the companion note
catalogues them~\cite{leannote}. Exact rational
arithmetic is likewise essential for the \emph{equality} deliverables,
decidable closure $t^\star=J^\star$ and exact contact $f_j(x^\star)=0$. Both
are $\Q$-decidable and auditable, but no finite-precision enclosure of
positive width can establish them: such an enclosure can bound a gap below
$\varepsilon$, never certify it is $0$. A zero-width enclosure can, and
fixed precision delivers one only in the exact-witness case named above. We
claim exactness for auditability and decidable equality, not as the only
route to a rigorous bound.

\section{Scope and limitations}\label{sec:limits}

\paragraph{Without A2, the deliverable is a bound, not closure.} Without A2,
closure is unattainable; the failure of A2 is generic when the active
keep-out guards a sub-block of dimension greater than one
(Remark~\ref{rem:contact}). The deliverable does not collapse: the sweep
still returns an exact, tolerance-free rational \emph{lower bound}, since
weak duality needs no rationality of the optimum. A keep-out guarding a
single coordinate with rational radius instead pins the active coordinate to
a rational value, leaving the remaining coordinates a rational linear solve.
Quantifying the resulting bracket (rates in the rationalization budget,
behaviour at the relaxation-gap floor) is outside this paper's scope.

\paragraph{Verification bit-complexity is measured and coarsely bounded, not
sharply bounded.} This paper \emph{measures} verification cost --- the linear
pivot growth of \S\ref{sec:rpod} --- and bounds it coarsely
(Remark~\ref{rem:effective}); sharp growth laws and their model sensitivity
are outside scope. That verification cost can grow along the horizon while
the emitted certificate stays bounded is visible in the demonstration of
\S\ref{sec:rpod}.

\paragraph{The demonstration family's mildness is tied to its shared
transition matrix.} The CW-coupled family couples every stage through one
shared rationalized transition matrix, computed for a circular reference
orbit. An eccentric-reference variant of the same construction --- the
Yamanaka--Ankersen transition matrix, a different matrix at every stage ---
reproduces the CW pivot sequence bit for bit at $e=0$. At the probed
eccentricities $e\in\{0.1,0.3,0.5\}$, rationalized to the same realized
fidelity, its sweep pivots
grow at $\approx1000$--$1050$ bits per stage against the CW family's
$\approx778$, and the driver is the per-stage time-variance of the
transition matrix rather than the eccentricity magnitude. On the CW family
the certified value's denominator is horizon-bounded (\S\ref{sec:rpod}); on
the variant the value's bit-size grows at $\approx1000$ bits per stage, so
the compactness of the recomputed value does not carry over. An executable
receipt catalogued in the companion note pins these
figures~\cite{leannote}.

\paragraph{The per-stage Hadamard/Cramer reconstruction bound is not
included.} A per-stage Hadamard/Cramer bound on the verifier's multiplier
reconstruction is weaker than the $2L$ bound the design order gives, a
consequence of the design order rather than a route to it, and not needed
for either leg of Theorem~\ref{thm:main}.

\paragraph{Bienstock realizability inside the class is not established.}
Quadratic optimization at large admits $2^n$-bit
optima~\cite{bienstock2020}; whether such a gadget can be realised as an
instance of~\eqref{eq:qcqp} is not settled either way. That is why
Theorem~\ref{thm:main} is stated conditionally, and why this paper does not
claim conditioning is unavoidable.

\paragraph{Hard equality dynamics are outside the class.} Hard equality
dynamics are outside~\eqref{eq:qcqp} by definition (\S\ref{sec:regime});
readers of the banded-KKT literature will recognise that Wright and
Rao--Wright--Rawlings~\cite{wright1993,rao1998} factorize precisely that
equality-constrained system, in floating point. The exact-certificate
treatment of hard dynamics itself is outside scope.

\paragraph{A1-failing keep-out geometry needs level two.} Under A1, multiple
keep-outs per stage already close (Theorem~\ref{thm:main}); when a stage's
keep-out geometry itself violates A1 --- several exclusion balls on one
stage, two quadratic constraints where the S-lemma wants one, the
sign-flipped analogue of the two-\emph{inclusion}-constraint CDT
subproblem~\cite{celis1985,aizhang2009} --- level-one certification gaps,
and an exact level-two treatment is outside scope.

\section{Conclusion}\label{sec:conclusion}
The exact-rational certification of a \emph{penalty-coupled} block-banded
trajectory QCQP is an LQ Riccati sweep run in $\Q$. The sweep is sound at the
band-Hessian test (Lemma~\ref{lem:band}); the border-carrying extension's
soundness is asserted and test-pinned, its count read off the sweep's loop
bounds (\S\ref{sec:riccati}); the trusted path is
machine-checked~\cite{leannote}; and at a constructed rational optimum the
bracket closes under an aggregate-PSD condition implied by a local
diagonal-dominance rule (Theorem~\ref{thm:main}, Lemma~\ref{lem:dd}).

Operationally, exactness buys \emph{auditability}: the certificate is a
finite rational object that \emph{any} independent implementation \emph{can}
re-verify bit-for-bit, with no tolerance to agree on and no dependence on
platform rounding. Three conditions govern that re-verification --- exact
rational arithmetic, unbounded integers, and the transmitted data; they are
the re-verifier's and not the certificate's, and each is load-bearing. The
same sweep transcribed operation for operation into IEEE-754 doubles
\emph{refuses} a published closing certificate --- the disjoint-contact
two-keep-out instance of \S\ref{sec:theorem} --- its terminal corner coming
out $-4.44\times10^{-16}$ rather than exactly $0$: a tolerance-free test has
no tolerance to absorb its own rounding. The $t^\star$ the verifier recomputes on the
demonstration family of \S\ref{sec:rpod} runs $931$--$1089$ bits, which no
fixed-width arithmetic can hold, while the emitted pair itself stays at
$6$--$8$ bits. And Definition~\ref{def:protocol} conditions the certificate
on agreed data, so a re-verifier that re-derives the rationalized $\Phi$ of
\S\ref{sec:rpod} instead of reading the transmitted rationals certifies a
different instance. So conditioned, the acceptance of a penalty-coupled
trajectory QCQP in the certified class becomes a reproducible theorem about
the deployed numbers rather than a solver verdict that must itself be
trusted.

\bibliographystyle{plain}
\bibliography{references}

\end{document}
